% % ######% ######% ######% ######% ######% ######% ######% ######% ######% ######
\section{Impact of Work Stealing}
In this experiment, we explore the degree of freedom for pulling work from other threads when a worker is idle. We want to answer the research question

\begin{tcolorbox}[colback=blue!5, colframe=blue!40!black, title=Research Questions]
    \textbf{RQ:} If a thread has no work, from which thread should it steal?
\end{tcolorbox}

% % ######% ######% ######% ######% ######% ######% ######% ######% ######% ######
\subsection{Variants}
Work stealing is the complement of work dealing: while dealing determines where to \emph{push} successor tasks, stealing determines where to \emph{pull} tasks when a worker's own queue is empty.
In our SPE, when a worker wakes up and finds its per-thread queue empty, it can attempt to steal a task from another worker's queue before falling back to the admission queue.
The stealing strategy determines how the victim queue is selected.

We fix the work-dealing strategy to \texttt{PRODUCER\_LOCAL} for all experiments in this section, as it was the best-performing dealing strategy on successor-heavy workloads.
This isolates the effect of stealing: successor tasks stay local by default, and stealing only occurs when a thread has genuinely run out of work.

\subsubsection{No Stealing}
\texttt{NONE:}
The worker skips stealing entirely and falls through directly to the admission queue.
This is the baseline---it preserves maximum locality but risks idle threads when load is imbalanced.

\subsubsection{Stealing Strategies}
The following strategies are invoked after a worker finds its own per-thread queue empty and before it falls back to the admission queue.

\texttt{ROUND\_ROBIN:}
The worker tries victims in order starting from its next neighbor:
\[
    \text{victim} = (i+1) \bmod N, \quad (i+2) \bmod N, \quad \ldots
\]
This scans all $N-1$ queues in the worst case.

\textit{Advantages:} Exhaustive search guarantees finding a task if one exists.

\textit{Disadvantages:} $O(N)$ scan cost. Always starts with the same neighbor, creating contention hotspots when multiple threads steal simultaneously. The scan itself contends on the victim queues' atomic dequeue operations.

\texttt{RANDOM:}
The worker picks one random victim queue using a thread-local RNG.

\textit{Advantages:} $O(1)$ cost. No deterministic contention patterns.

\textit{Disadvantages:} May miss tasks (only one attempt). The randomly chosen victim may be empty while other queues have work.

\texttt{CHOOSE\_TWO:}
The worker samples two random queues and steals from the one with more pending tasks:
\[
    a, b \sim \text{Uniform}(0, N), \quad \text{victim} =
    \begin{cases}
        a & \text{if } \texttt{approxSize}[a] \geq \texttt{approxSize}[b] \\
        b & \text{otherwise}
    \end{cases}
\]
If the first choice fails, it tries the second.

\textit{Advantages:} $O(1)$ cost with a better chance of finding a loaded queue than \texttt{RANDOM}. Two attempts provide a fallback.

\textit{Disadvantages:} Two atomic reads for size comparison. May still miss if both sampled queues are empty.

\texttt{LARGEST\_QUEUE:}
The worker scans all queues and steals from the one with the most pending tasks:
\[
    \text{victim} = \arg\max_{j \neq i} \texttt{approxSize}[j]
\]

\textit{Advantages:} Steals from the most overloaded thread, providing the best load balancing per steal.

\textit{Disadvantages:} $O(N)$ scan cost, same as \texttt{ROUND\_ROBIN}. The scan reads all approximate size counters, introducing cache line traffic even when no task is ultimately stolen.

% % ######% ######% ######% ######% ######% ######% ######% ######% ######% ######
\subsection{Experimental Results: Work Stealing Strategies}
\label{sec:stealing-results}

We evaluate the five work-stealing strategies on the same four workloads as in the dealing experiments, with dealing fixed to \texttt{PRODUCER\_LOCAL}.

\subsubsection{Stateless Workload}

Table~\ref{tab:ws-stateless} shows the results. Since the stateless query produces no successor tasks, the per-thread queues are always empty and stealing has no meaningful work to find.

\begin{table}[h]
\centering
\caption{Stateless workload throughput (MB/s) by stealing strategy (dealing = \texttt{PRODUCER\_LOCAL}).}
\label{tab:ws-stateless}
\begin{tabular}{l|rrrrrrr}
\hline
\textbf{Stealing} & \textbf{1T} & \textbf{2T} & \textbf{4T} & \textbf{8T} & \textbf{16T} & \textbf{32T} & \textbf{64T} \\
\hline
\texttt{NONE}          & 78 & \textbf{152} & \textbf{300} & \textbf{586} & 593 & 661 & 507 \\
\texttt{ROUND\_ROBIN}  & 78 & 150 & 294 & 564 & 571 & 679 & \textbf{524} \\
\texttt{RANDOM}        & 77 & 151 & 298 & 584 & 591 & 661 & 499 \\
\texttt{CHOOSE\_TWO}   & 77 & 151 & 292 & 578 & 592 & 663 & 490 \\
\texttt{LARGEST\_QUEUE}& 77 & 151 & 298 & 581 & \textbf{593} & \textbf{702} & 510 \\
\hline
\end{tabular}
\end{table}

The key observations are as follows.
%
\begin{itemize}
    \item \textbf{All strategies perform within 5\%} at every thread count, confirming that stealing has negligible impact when there are no successor tasks to steal.
    \item \textbf{\texttt{NONE} is marginally best at 2--8T} because the stealing check itself introduces a small overhead (failed dequeue attempts on empty queues).
    \item \textbf{\texttt{LARGEST\_QUEUE} leads at 32T} (702\,MB/s), possibly because its full scan of approximate sizes acts as a brief pause that reduces contention on the admission queue.
\end{itemize}

\subsubsection{High-Frequency Window Workload}

Table~\ref{tab:ws-highfreq} shows the results. With $\sim$5 million successor tasks, stealing strategies can now find work in other threads' queues.

\begin{table}[h]
\centering
\caption{High-frequency window workload throughput (MB/s) by stealing strategy (dealing = \texttt{PRODUCER\_LOCAL}).}
\label{tab:ws-highfreq}
\begin{tabular}{l|rrrrrrr}
\hline
\textbf{Stealing} & \textbf{1T} & \textbf{2T} & \textbf{4T} & \textbf{8T} & \textbf{16T} & \textbf{32T} & \textbf{64T} \\
\hline
\texttt{NONE}          & \textbf{39} & \textbf{52} & \textbf{87} & \textbf{125} & \textbf{130} & 128 & 39 \\
\texttt{ROUND\_ROBIN}  & 39 & 34 & 32 & 32 & 32 & 31 & 29 \\
\texttt{RANDOM}        & 39 & 50 & 81 & 118 & 121 & \textbf{131} & \textbf{42} \\
\texttt{CHOOSE\_TWO}   & 39 & 47 & 73 & 108 & 110 & 130 & \textbf{46} \\
\texttt{LARGEST\_QUEUE}& 38 & 34 & 32 & 32 & 32 & 31 & 29 \\
\hline
\end{tabular}
\end{table}

The key observations are as follows.
%
\begin{itemize}
    \item \textbf{\texttt{ROUND\_ROBIN} and \texttt{LARGEST\_QUEUE} are catastrophic}, dropping throughput from 125 to 32\,MB/s at 8T ($-$75\%).
    Both strategies scan all $N$ queues on every steal attempt. When combined with \texttt{PRODUCER\_LOCAL} dealing (where successor tasks are local), the scan contends with the producing thread's enqueue operations, effectively serializing the system.
    The contention is proportional to the number of threads, explaining why the throughput flatlines at $\sim$32\,MB/s regardless of thread count.

    \item \textbf{\texttt{NONE} is the best strategy at 1--16T}, confirming that when dealing is \texttt{PRODUCER\_LOCAL}, stealing is counterproductive.
    The producer-local policy already ensures tasks are on the right thread; stealing disrupts this by moving tasks to threads where the data is cache-cold.

    \item \textbf{\texttt{RANDOM} is the safest non-\texttt{NONE} option} (118\,MB/s at 8T vs.\ 125 for \texttt{NONE}, only $-$6\%).
    Its $O(1)$ single-attempt approach minimizes contention.
    At 32--64T, it slightly outperforms \texttt{NONE} (131/42 vs.\ 128/39), suggesting that at high thread counts, occasional stealing helps rebalance across hyperthreads.

    \item \textbf{\texttt{CHOOSE\_TWO} trails \texttt{RANDOM}} (108\,MB/s at 8T) because it reads two approximate sizes before stealing, doubling the cache line traffic compared to a single random attempt.
\end{itemize}

\subsubsection{Join Workload}

Table~\ref{tab:ws-join} shows the results.

\begin{table}[h]
\centering
\caption{Join workload throughput (MB/s) by stealing strategy (dealing = \texttt{PRODUCER\_LOCAL}).}
\label{tab:ws-join}
\begin{tabular}{l|rrrrrrr}
\hline
\textbf{Stealing} & \textbf{1T} & \textbf{2T} & \textbf{4T} & \textbf{8T} & \textbf{16T} & \textbf{32T} & \textbf{64T} \\
\hline
\texttt{NONE}          & 31 & 41 & \textbf{65} & \textbf{100} & 109 & 87 & 61 \\
\texttt{ROUND\_ROBIN}  & 31 & 41 & 62 & 89 & 100 & 104 & 65 \\
\texttt{RANDOM}        & 31 & \textbf{42} & 61 & 90 & 105 & 96 & 62 \\
\texttt{CHOOSE\_TWO}   & 31 & \textbf{42} & 62 & 87 & 101 & \textbf{110} & \textbf{70} \\
\texttt{LARGEST\_QUEUE}& 31 & 41 & 62 & 90 & 99 & 103 & 65 \\
\hline
\end{tabular}
\end{table}

The key observations are as follows.
%
\begin{itemize}
    \item \textbf{The join workload shows a crossover: \texttt{NONE} leads at 4--8T, but \texttt{CHOOSE\_TWO} wins at 32--64T.}
    At low thread counts, locality dominates (same as highfreq). At high thread counts, the join creates load imbalance between the three pipelines (left source, right source, join), which stealing can address.
    \texttt{CHOOSE\_TWO} achieves 110\,MB/s at 32T (+26\% over \texttt{NONE}'s 87\,MB/s) and 70 at 64T (+15\%).

    \item \textbf{\texttt{ROUND\_ROBIN} stealing is less harmful here than on highfreq} (89\,MB/s at 8T vs.\ 32\,MB/s on highfreq).
    The join pipeline has fewer but larger successor tasks (join result buffers), so the scan contention is amortized over more computation per task.

    \item \textbf{The join is the only workload where stealing provides a net benefit} at any thread count, because the multi-pipeline structure inherently creates imbalance that dealing alone cannot resolve.
\end{itemize}

\subsubsection{Mixed Concurrent Workload}

Table~\ref{tab:ws-mixed} shows wall-clock time (lower is better).

\begin{table}[h]
\centering
\caption{Mixed workload wall-clock time (seconds, lower is better) by stealing strategy (dealing = \texttt{PRODUCER\_LOCAL}).}
\label{tab:ws-mixed}
\begin{tabular}{l|rrrrrrr}
\hline
\textbf{Stealing} & \textbf{1T} & \textbf{2T} & \textbf{4T} & \textbf{8T} & \textbf{16T} & \textbf{32T} & \textbf{64T} \\
\hline
\texttt{NONE}          & 198 & 125 & \textbf{71} & \textbf{45} & \textbf{42} & \textbf{40} & 112 \\
\texttt{ROUND\_ROBIN}  & 196 & 127 & 85 & 67 & 92 & 117 & 148 \\
\texttt{RANDOM}        & 196 & \textbf{124} & 73 & 48 & 45 & \textbf{40} & \textbf{109} \\
\texttt{CHOOSE\_TWO}   & 196 & 125 & 76 & 52 & 48 & 41 & 103 \\
\texttt{LARGEST\_QUEUE}& 198 & 126 & 85 & 68 & 92 & 117 & 144 \\
\hline
\end{tabular}
\end{table}

The key observations are as follows.
%
\begin{itemize}
    \item \textbf{\texttt{NONE} is the best at 4--32T}, consistent with the highfreq results.
    The mixed workload's asymmetric load (fast filter + slow aggregation) is better handled by dealing (\texttt{PRODUCER\_LOCAL} keeps aggregation successors local) than by stealing.

    \item \textbf{\texttt{ROUND\_ROBIN} and \texttt{LARGEST\_QUEUE} are again catastrophic} at 16--32T (92--117\,s vs.\ 42--40\,s for \texttt{NONE}).
    The $O(N)$ scan during stealing contends with the aggregation pipeline's high-frequency successor task production, creating a 2--3$\times$ slowdown.

    \item \textbf{\texttt{CHOOSE\_TWO} is the best stealing strategy at 64T} (103\,s vs.\ 112 for \texttt{NONE}, $-$8\%).
    Cross-NUMA execution creates sufficient imbalance that targeted stealing (from the fuller of two random queues) provides a net benefit despite the locality cost.

    \item \textbf{\texttt{RANDOM} is a reliable middle ground} across all thread counts: never more than 7\% worse than \texttt{NONE}, and occasionally better at 64T.
\end{itemize}

\subsection{Summary}

\begin{itemize}
    \item \textbf{\texttt{NONE} (no stealing) is the best default} when combined with \texttt{PRODUCER\_LOCAL} dealing.
    Stealing disrupts the locality that \texttt{PRODUCER\_LOCAL} establishes; the overhead of failed steal attempts and the cache pollution from stolen tasks outweigh any load-balancing benefit.

    \item \textbf{\texttt{ROUND\_ROBIN} and \texttt{LARGEST\_QUEUE} stealing should be avoided.}
    Their $O(N)$ scan creates devastating contention on successor-heavy workloads, causing up to 75\% throughput regression on the high-frequency window workload and 2--3$\times$ slowdown on the mixed workload.

    \item \textbf{\texttt{CHOOSE\_TWO} stealing is beneficial only on multi-pipeline workloads at $\geq 32$T.}
    The join workload is the sole scenario where stealing provides a net benefit (+26\% at 32T), because the multi-pipeline structure creates inherent load imbalance that dealing alone cannot resolve.

    \item \textbf{\texttt{RANDOM} is the safest non-\texttt{NONE} option:} $O(1)$ cost, minimal contention, never more than 6\% worse than \texttt{NONE}, and occasionally better at high thread counts.

    \item \textbf{Work stealing and work dealing are complementary but not additive.}
    The best overall configuration is \texttt{PRODUCER\_LOCAL} dealing with \texttt{NONE} stealing (locality-first).
    Enabling stealing on top of \texttt{PRODUCER\_LOCAL} generally hurts, because the tasks that would be stolen are exactly the ones that benefit from staying local.
    Stealing is most justified when dealing distributes tasks evenly (e.g., \texttt{ROUND\_ROBIN}) and load imbalance arises from uneven task execution times rather than uneven dispatch.
\end{itemize}

\subsection{Discussion}
